% TOP_PROBLEM: {"id": 4700006, "title": "Trigonometric Abel differential equations I", "category_id": 11, "category": {"id": 11, "name": "dynamical_systems", "display_name": "Dynamical Systems", "description": "Problems about long-term behavior of deterministic systems, Hamiltonian dynamics, and periodic orbits.", "slug": "dynamical-systems", "order_index": 11, "created_at": "2026-07-31T15:26:25.671Z"}, "status": "open", "problem_number": "AMR-046-0006", "set_id": 13}
% FIRST_POSTED: 2026-10-02
% ENTRY_KIND: statement_only
% UnsolvedMath ID 4700006; source catalogue code AMR-046-0006
% !TeX program = lualatex
% Definition and sources only; this file is not an LLM solution attempt.
\documentclass[11pt,a4paper]{article}
\usepackage[margin=25mm]{geometry}
\usepackage{fontspec}
\setmainfont{lmroman10-regular.otf}[BoldFont=lmroman10-bold.otf,
 ItalicFont=lmroman10-italic.otf,BoldItalicFont=lmroman10-bolditalic.otf]
\usepackage{amsmath,amssymb,mathtools}
\usepackage{unicode-math}
\setmathfont{latinmodern-math.otf}
\AtBeginDocument{\let\setminus\smallsetminus}
\usepackage{xurl,hyperref}
\hypersetup{colorlinks=true,urlcolor=blue}
\setcounter{secnumdepth}{0}
\setlength{\parindent}{0pt}
\setlength{\parskip}{5pt}
\setlength{\emergencystretch}{3em}
\begin{document}
\section*{Trigonometric Abel differential equations I}\label{4700006}
{\small UnsolvedMath ID 4700006 · AMR-046-0006 · Dynamical Systems · AMR Open Problem Lists}
\subsection{Definitions and mathematical statement}
For arbitrary real $a_0,a_1,a_2,b_0,b_1,b_2$, consider the scalar equation
\[
 x'=(a_0+a_1\sin t+a_2\cos t)x^3
       +(b_0+b_1\sin t+b_2\cos t)x^2.
\]
Only real solutions defined and finite on an entire interval of length
$2\pi$ are considered. Such a solution is periodic if it extends by
$x(t+2\pi)=x(t)$ to all real times. Its graph on the cylinder
$(\mathbb R/2\pi\mathbb Z)\times\mathbb R$ is a limit cycle when its initial
value is an isolated fixed point of the time-$2\pi$ return map.

The question is whether the maximum number of distinct limit cycles, as the
six coefficients vary, equals three. The count includes the zero solution
when it is isolated, as well as both positive and negative periodic
solutions. No multiplicities of return-map zeros are included in this count.
Solutions with a pole during the period are excluded, even if a formal
projective continuation would return to the same value.

Because $x=0$ is a solution and uniqueness prevents another solution from
crossing it, every nonzero periodic solution has a fixed sign. This separates
the two half-cylinders but does not remove either from the question. If the
return map is the identity on an interval, its continuum of fixed points
does not consist of isolated limit cycles. The requested bound must also
cover coefficients that change sign and all such degenerate parameter cases.
\subsection{Short English statement}
Can a first-harmonic trigonometric Abel equation have more than three isolated periodic solutions, counting zero when isolated?
\subsection{Sources}
\begin{itemize}
\item[S1] ULAM.AI and the UnsolvedMath Contributors, supplied problems.json, record 4700006 (AMR-046-0006), AMR Open Problem Lists. Catalogue identity and source transcription; CC BY 4.0.
\url{https://huggingface.co/datasets/ulamai/UnsolvedMath}.
\item[S2] Gasull - Some open problems in low dimensional dynamical systems (2020), Problem environment 6: Trigonometric Abel differential equations I. Original source indexed by AMR-046-0006.
\url{https://arxiv.org/abs/2012.02524}.
\end{itemize}
\end{document}
